\section{按辐射度量域分阶段优化}
\label{sec:optimization}

\subsection{损失在哪个域计算}
\label{sec:lossdomain}
重光照的光度损失通常在显示编码图像上计算，这也是 LDR 拍摄数据的原生格式。
我们使用纯幂律 $\Gamma(L)=L^{1/\gamma}$、$\gamma=2.2$ 显示 HDR 数据，并对 LDR 数据做相应的逆变换以线性化；此时，$\ell_1$ 损失对预测辐亮度 $\hat L$ 的梯度幅值为
\begin{equation}
  \Gamma'(\hat L)=\tfrac{1}{\gamma}\,\hat L^{\,1/\gamma-1},
  \label{eq:slope}
\end{equation}
与残差大小无关。
因此，即使残差很小，预测值仅为白色值百分之一的像素，获得的梯度也会是白色像素的十二倍。
由于 \cref{eq:radiance} 的预测是多个可学习因子的乘积，训练初期像素偏暗，往往只是因为可见性、法线或材质尚不准确。
我们推测，这些像素因而主导了传向几何的梯度，优化器则通过平滑细小结构来降低损失。
\cref{sec:exp_optimization} 的受控实验表明，在我们的模型中，损失计算域决定了细小几何能否形成；但实验并未单独识别其中机制，我们也尚未测试在 $\Gamma$ 内加入偏移量或对损失重加权等补救方式~\cite{mildenhall2022nerf}。

\subsection{三个阶段}
\label{sec:stages}
我们在 32 张训练轮廓构成的视觉外壳内均匀初始化 20k 个高斯（一个点落在至少 90\% 的轮廓内才被保留），并将初始法线设为背离场景中心的方向。

\paragraph{I. 几何预热（8k 次迭代）。}
不含网络的点光源模型以 $\softplus(\vect b_i)\,\bar E(\mu_i)\,\psi(\vect n_i\cdot\vect l_i)$ 为每个高斯赋色，仅优化几何、基础颜色和法线。
渲染器始终输出线性值，而目标随训练变化：预热的前三分之二使用显示编码图像 $\Gamma(Y)$，之后将其指数从 $1/\gamma$ 线性退火至 $1$。
因此，几何先在动态范围被压缩的目标下形成，再平滑过渡至线性辐亮度。
\gsthree 和 SSD-GS 对 HDR 训练图像采用了类似退火~\cite{bi2024gs3,zheng2026ssdgs}；我们将其用于所有数据类别，并使后续联合训练保持在线性辐亮度域。

\paragraph{II. 联合训练（22k 次迭代）。}
以高斯为接收点，使用线性辐亮度损失训练 \cref{eq:radiance} 的完整模型；增密和剪枝遵循 3DGS~\cite{kerbl20233d}，持续至第 25k 次迭代，高斯数量最多为 400k。

\paragraph{III. 外观精修（8k 次迭代）。}
冻结几何和相机标定，将接收点切换为像素，并将损失移回图像评估所使用的显示域。
实验表明，几何固定后，对暗像素施加的额外权重能使两种接收点的三项指标均得到提升（\cref{sec:exp_optimization}）。

所有阶段均最小化 $0.8\,\ell_1+0.2\,(1-\mathrm{SSIM})$，两项都在当前目标域计算；此外，利用基准提供的前景掩码 $\alpha$，对渲染不透明度加入 $0.5\,\|\hat\alpha-\alpha\|_1$，并对编码 $\vect f$ 施加较弱的 $\ell_2$ 惩罚。

\subsection{实现细节}
\label{sec:implementation}
\ours 使用 PyTorch 和 gsplat~\cite{ye2025gsplat}，图集分辨率为 $512^2$ 纹素，具有 $K=7$ 个层级和 $C=8$ 个通量通道。
MLP 宽度为 128（可见性网络为 64），初始化使训练从 $L_{\mathrm{tr}}\approx 0$ 和 $V=V_0$ 开始；采用 Adam 优化器（网络学习率为 $10^{-3}$，衰减至约四分之一；高斯参数沿用 3DGS）。
对真实拍摄数据，我们还优化每个训练相机绕其标定中心的旋转，并在每一步后去除这类旋转调整可以吸收的全局平移，使重建保持在标定坐标系内；测试相机和光源始终不作调整。
全部 18 个场景使用同一配置，仅按各基准的数据约定做适配；在单张 NVIDIA RTX 6000 Ada 上，一次完整训练耗时 17 至 28 分钟，平均为 21 分钟。
